\documentclass[11pt]{article}
\usepackage[margin=1in]{geometry}
\usepackage{amsmath,amssymb,amsthm,hyperref}
\setlength{\emergencystretch}{2em}
\newtheorem{lemma}{Lemma}
\newtheorem{theorem}{Theorem}
\title{Equal perturbation norms do not characterize\\maximal monotonicity: Conjecture 00000008852}
\author{ziangni-sys}
\date{October 8, 2026}
\begin{document}
\maketitle

\paragraph{The clause being disproved.}
The English statement says that the perturbed operator remains maximal
monotone \emph{exactly when} its perturbation norm $\delta$ does not exceed
the strong monotonicity constant $\gamma$ of the original operator.
We refute this necessary-and-sufficient claim for individual perturbations,
even for bounded linear operators on the real Hilbert space $\mathbb R$,
using the optimal strong monotonicity constant. The counterexamples can
have arbitrarily small absolute perturbation norm.
The Chinese version describes a threshold at $\delta\leq\gamma$.
A \emph{uniform worst-case guarantee} over all perturbations of bounded norm
is a different statement: our example does not refute that sufficient
guarantee or a sharp worst-case threshold. It refutes the pointwise
``exactly when'' assertion in the supplied English statement.
Since this is a conjunct of the conjecture, the other trajectory assertions
need not be decided.

\paragraph{Definitions.}
On $H=\mathbb R$, the inner product is $\langle x,y\rangle=xy$ and the norm is
$|x|$. An operator $F$ is $\gamma$-strongly monotone when
\[
 (x-y)(F(x)-F(y))\geq\gamma |x-y|^2 \qquad(x,y\in\mathbb R).
\]
A relation $R\subset\mathbb R^2$ is monotone if
$(x-y)(u-v)\geq0$ whenever $(x,u),(y,v)\in R$.
An operator is maximal monotone if its graph is monotone and has no proper
monotone relational extension. Thus maximality here is not merely maximality
among bounded linear operators.

\begin{lemma}
For every $t>0$, the operator $F_t(x)=tx$ is maximal monotone.
\end{lemma}
\begin{proof}
Graph monotonicity is $t(x-y)^2\geq0$. Suppose a monotone relation $R$
contains the graph, and $(x,u)\in R$. Put $y=(u+tx)/(2t)$.
The point $(y,ty)$ belongs to the graph, hence to $R$, so
\[
 0\leq(x-y)(u-ty)=-\frac{(u-tx)^2}{4t}\leq0.
\]
It follows that $u=tx$. Every point of $R$ already belongs to the graph,
which proves full relational maximality.
\end{proof}

\begin{theorem}
For any $a>0$, let $A=aI$ and $B_+=2aI$, $B_-=-2aI$. Then:
\begin{enumerate}
\item The optimal strong monotonicity constant of $A$ is $\gamma=a$.
\item Both perturbations have operator norm $\delta=2a>\gamma$.
\item $A+B_+=3aI$ is maximal monotone.
\item $A+B_-=-aI$ is not even monotone.
\end{enumerate}
\end{theorem}
\begin{proof}
For $A$, the strong monotonicity inequality holds with equality at
$\gamma=a$. Conversely, testing that inequality at $x=1$, $y=0$ forces
every admissible constant to be at most $a$.
The bounded linear scalar operator $cI$ has norm $|c|$: its action has
$|cx|=|c||x|$ for all $x$, and the upper bound is attained at $x=1$.
This gives both stated perturbation norms. The preceding lemma with
$t=3a$ proves the positive case. In the negative case, graph points
$(1,-a)$ and $(0,0)$ have monotonicity pairing $-a<0$.
\end{proof}

\paragraph{Arbitrarily small perturbations and zeros.}
Given any tolerance $\varepsilon>0$, take $a=\varepsilon/4$.
Then $\|B_+\|=\varepsilon/2<\varepsilon$, but
$\|B_+\|>\gamma=\varepsilon/4$, while $A+B_+$ is maximal monotone.
This prevents an unspecified requirement of small \emph{absolute} norm from
saving the necessity claim. The unperturbed and both perturbed operators
have the unique zero $0$, so the example lies within the setting of operators
with a zero described by the source. In particular, zero existence causes
no difficulty. The opposite outcomes of $B_+$ and $B_-$ show concretely
why perturbation size alone cannot decide the property for a given operator.

\paragraph{Formalization.}
The Lean project uses Mathlib's actual continuous linear maps
$\mathbb R\to_L\mathbb R$ and their operator norm. It defines monotone
relations and maximality by quantification over all relations on the real
Hilbert space, and proves the maximality lemma directly.
The theorem \texttt{same\_norm\_opposite\_behavior} establishes the optimal
constant, both exact operator norms, and the opposite monotonicity outcomes.
The theorem \texttt{arbitrarily\_small\_counterexample} proves the tolerance
statement for every $\varepsilon>0$, and \texttt{conjectured\_iff\_false}
negates the claimed criterion even when the constant is required to be
optimal. No graph maximality or spectral fact is assumed as a certificate.
The final theorem concerns the disputed conjunct, not a replacement
statement about convergence or continuity of zeros.
\end{document}
